\documentclass[Total.tex]{subfiles}

\begin{document}
	\chapter{Homotopy Groups}
		Recall:
		\begin{itemize}
			\item Component $\pi_0$;
			\item Fundamental Group $\pi_1$;
		\end{itemize}
		
		And with covering $p:X\rightarrow Y$, we obtain the induced injection
		\begin{gather*}
			\pi_1(p,x):\pi_1(X,x)\rightarrow \pi_1(Y,y)
		\end{gather*}
		
		
		Especially, for a $G$-covering, an exact sequence
		\begin{gather*}
			1\rightarrow \pi_1(X,x)\stackrel{\pi_1(p,x)}{\rightarrow} \pi_1(Y,y) \stackrel{\partial}{\rightarrow} G\rightarrow 1
		\end{gather*};
		
		Seifert-van Kampen: 
		
		$Aut_{\prod(X)}(x)$;
		
		Homotopy $f_0,f_1:X\rightarrow Y$, then we obtain a group isomorphism $\pi_1(Y,y_0)\cong \pi_1(Y,y_1)$:
		\begin{center}
			\begin{tikzcd}
				&\pi_1(Y,y_0)\arrow[dd]\\
				\pi_1(X,x)\arrow[ru]\arrow[rd]\\
				&\pi_1(Y,y_1)
			\end{tikzcd}
		\end{center}
		
		\subsubsection{LES}
		
		Consider the following triple $(X,A,x)$. The \textbf{long exact sequence}
		\begin{gather*}
			\pi_1(A,x)\stackrel{\pi_1(i,x)}{\rightarrow} 
		\end{gather*}
		\section{Homotopy Group}
			\begin{Def}
				Define:
				\begin{itemize}
					\item $I^n$; $\partial I^n=\{(s_1,\dots,s_n)|\exists i,s_i\in \{0,1\}\}$
					\item Homotopy Classes $\pi_n(X,x)=\{[f]|f:(I^n,\partial I^n)\rightarrow (X,\{x\})\}$;
				\end{itemize}
				
				Then: Group Structure
				\begin{itemize}
					\item Multiplication:
					\begin{gather*}
						f*g:=\begin{cases}
							f(2s_1,s_2,\dots,s_n)&\\
							f(2s_1-1,s_2,\dots,s_n)
						\end{cases}
					\end{gather*}
					\item Identity:$s\mapsto x$;
					\item Inverse: $f^-(s_1,\dots,s_n):=(1-s_1,s_2,\dots,s_n)$;
				\end{itemize}
			\end{Def}
			
			\begin{proposition}
				The group $\pi_n(X,x)$ is abelian, $n\geq 2$. 
			\end{proposition}
			\begin{proof}
				We observe, that
			\end{proof}
			\subsubsection{Definition through Loop Space}
			\subsubsection{Factorization through Sphere}
			\subsection{Functorial Property}
			\subsection{Coverings}
			
			\begin{theorem}
				(Homotopy Groups and Coverings) Let $p:X\rightarrow Y$ be a covering, choose a base point $x\in X$ with $p(x)=y$, then
				\begin{gather*}
					内容...
				\end{gather*}
			\end{theorem}
			\subsection{Homotopy Groups of Pairs and Triple,Long Exact Sequence}
		\section{$n$-Connectivity}
			\begin{Def}
				\begin{itemize}
					\item A space $X$ is called \textbf{$0$-connected} if $\pi_0(X)$ consists of one point; (Path-connected)
					\item A space $X$ is called \textbf{$1$-connected} if $\pi_k(X)=\bullet,\forall 1\leq k\leq n$;
					\item A space $X$ is called \textbf{$\infty$-connected} if $\pi_k(X)=\bullet,k\geq 1$;
				\end{itemize}
			\end{Def}
			
			And:
			\begin{Def} Let $f:X\rightarrow X$;
				\begin{itemize}
					\item It is called $0$-connected if the induced map 
					\begin{gather*}
						\pi_0(f):\pi_0(X,x)\rightarrow \pi_0(X,x)
					\end{gather*}
					
					is surjective;
					\item It is called \textbf{$n$-connected} of the induced map
					\begin{gather*}
						\pi_k(f):\pi_k(X,x)\rightarrow \pi_k(X,x)\qquad\mbox{is bijective},0\leq k\leq n-1\\
						\pi_n(f):\pi_n(X,x)\rightarrow \pi_n(X,x)\qquad\mbox{is surjective}
					\end{gather*}
					\item It is called \textbf{$\infty$-connected} if $\pi_k(f):\pi_k(X,x)\rightarrow \pi_k(X,x)$ is bijective;
				\end{itemize}
			\end{Def}
			
			\begin{proposition}
				(Criteria by Sphere) TFAE:
				\begin{itemize}
					\item $\pi_n(X,x)$ is trivial for any base point $x\in X$;
					\item Every $S^n\rightarrow X$ is nullhomotopic. i.e. $[S^n,X]=id_\bullet$;
					\item Every $S^n\rightarrow X$ extends to $D^{n+1}\rightarrow X$;
				\end{itemize}
			\end{proposition}
			\begin{proof}
				内容...
			\end{proof}
			
			\begin{corollary}
				(Criteria by Sphere) TFAE:
				\begin{itemize}
					\item $X$ is $n$-connected;
					\item Given any $k$ with $0\leq k\leq n$, every map $S^k\rightarrow X$ is nullhomotopic;
					\item Given any $k$ with $0\leq k\leq n$, every map $S^k\rightarrow X$ extends to a map $D^{k+1}\rightarrow X$;
				\end{itemize}
			\end{corollary}
			\begin{proof}
				Omitted. 
			\end{proof}
			
			Moreover, 
			\begin{corollary}
				内容...
			\end{corollary}
			\begin{proof}
				内容...
			\end{proof}
			
		\section{Colimits}
			\begin{theorem}
				(Homotopy Groups and Colimits) Let: $X$ be a topological Hausdorff space, with a sequence of \underline{closed subspaces}
				\begin{gather*}
					内容...
				\end{gather*}
			\end{theorem}
			\begin{proof}
				内容...
			\end{proof}
	\chapter{CW-Theory}
		
		\section{CW-Approxiamation}
		\section{Whitehead's Theorem}
	\chapter{Homotopy and Homology}
		\section{Hurewicz's Theorem}
	\chapter{Fibration and Cofibration}
		\section{Fibrations}
			\subsection{Homotopy Lifting Property}
			\subsection{Maps into Fibration}
	\chapter{Spectra}
		Preparations: Recall some operations of spaces
		\begin{itemize}
			\item $\Sigma X:=X\wedge S^1$ Reduced suspension;
			\item $\Omega X:=\{\gamma:I\rightarrow X|\gamma(0)=\gamma(1)\}$;
		\end{itemize}
			\subsection{Stable Equivalence}
		\section{Spectra}
			
			\subsection{Some of Spectrum}
				Because of adjunction
				\begin{gather*}
					\Sigma \dashv \Omega
				\end{gather*}
				
				The spectrum has $\Sigma$-Representation and $\Omega$-Representation
				
				
				\subsubsection{Pre-Spectrum}
					\begin{Def}
						A \textbf{prespectrum} or \textbf{sequential spectrum} is:
						\begin{itemize}
							\item $E=((E_n),(\sigma_n:\Sigma E_n\rightarrow E_{n+1}))$; or by adjunction
							\item $E=((E_n),(\sigma_n:E_n\rightarrow \Omega E_{n+1}))$;
						\end{itemize}
						
						The only requirement is: The structure maps are maps between pointed spaces;
						
						Such $E_n$ is called $n$-level set.
					\end{Def}
					
					Then, a \textbf{morphism} between the is a family of morphisms $(f_n:E_n\rightarrow E_n')$ such that
					\begin{gather*}
						f_{n+1}\circ\sigma_n=\sigma_n\circ (\Sigma f_n)
					\end{gather*}
					
					\begin{center}
						\begin{tikzcd}
							\dots&X_n\arrow[dd,bend right=40,"f_n",swap]\arrow[d]&X_{n+1}\arrow[d]\arrow[dd,bend right=40,"f_{n+1}",swap]\\
							\dots&\Sigma X_n\arrow[ru]\arrow[dd,bend left=40,"\Sigma f_n"]&\Sigma X_{n+1}\arrow[dd,bend left=40,"\Sigma f_{n+1}"]\\
							\dots&Y_n\arrow[d]&Y_{n+1}\arrow[d]\\
							\dots&\Sigma Y_n\arrow[ru]&\Sigma Y_{n+1}
						\end{tikzcd}
					\end{center}
					
					\begin{Def}
						Category of sequential spectrum $\mathcal{Sp}^\NN$;
					\end{Def}
					
					And
					\begin{Def}
						Let $X$ be a spectrum, define
						\begin{gather*}
							内容...
						\end{gather*}
					\end{Def}
					
					
					\begin{example}
						(Sphere Spectral)
					\end{example}
					
					\begin{theorem}
						内容...
					\end{theorem}
					
				\subsubsection{Stable Homotopy Groups}
					In principle, we should not introduce this concept in unstable homotopy theory... But the fact is, we must introduce it here
					
					\begin{Def}
						
					\end{Def}
				\subsubsection{Spectrum}
				\begin{Def}
					A \textbf{$\Omega$-spectrum} is:
					\begin{itemize}
						\item $((E_n),(\sigma_n:\Sigma E_n\rightarrow E_{n+1}))$, and $\sigma_n$ is a weak equivalence; or
						\item $((E_n),(\tilde{\sigma}_n:E_n\rightarrow \Omega E_{n+1}))$ and $\tilde{\sigma}$ is a weak equivalence;
					\end{itemize}
				\end{Def}
				
				\begin{example}
					(Eilenberg-Maclane Spectrum)
				\end{example}
				
			\subsection{Spectrification}
				\begin{example}
					(Sphere Spectrum) Consider following spectrum
					\begin{gather*}
						内容...
					\end{gather*}
				\end{example}
			
			\subsection{Categoricial Property}
			
		\section{Stability}
	\chapter{Excision Theorem, and Freudenthal Suspension Theorem}
	\chapter{Appendix: Serre's Finiteness Theorem for Spheres}
		\begin{theorem}
			(Serre) 
			\begin{gather*}
				\pi_k(S^n)=\begin{cases}
					\ZZ&n=0\\
					a
				\end{cases}
			\end{gather*}
		\end{theorem}
		
		\section{Nilpotence Theorem}
		
		\section{Computing of $\pi_*(S^n)$	}
\end{document}